%%
%% fall-lecture18.tex
%% 
%% Made by Alex Nelson <pqnelson@gmail.com>
%% Login   <alex@lisp>
%% 
%% Started on  2026-11-08T08:04:00-0800
%% Last update 2026-11-08T08:04:00-0800
%% 

\lecture{}

Today, we will conclude our detour into Thermodynamics, and then start
discussing ``quasi one dimensional flows'' (which will probably be
studied for the remainder of the quarter).

First, we finish obtaining quantities for ideal and perfect gases.

\begin{node}[Entropy of ideal gas]
For a general gas, we can write entropy as a function of two [state]
variables. Let us write $S=S(T,V)$. Then the difference equation
\begin{subequations}
  \begin{align}
\D S &= \left(\frac{\partial S}{\partial T}\right)_{V}\D T+ \left(\frac{\partial S}{\partial V}\right)_{T}\D V\\
&= \frac{c_{V}}{T}\D T+ \left(\frac{\partial S}{\partial V}\right)_{T}\D V\quad\mbox{by definition of $c_{V}$}\\
&= \frac{c_{V}}{T}\D T + \frac{R}{V}\D V
  \end{align}
\end{subequations}
where, in the last line, we use the Maxwell relation for Helmholtz free energy Eq~\eqref{eq:maxwell-relation-for-helmholtz-free-energy}
\begin{subequations}
  \begin{align}
\left(\frac{\partial S}{\partial V}\right)_{T}
&=\left(\frac{\partial p}{\partial T}\right)_{V}\\
&=\left(\frac{\partial RT/V}{\partial T}\right)_{V}=\frac{R}{V}.
  \end{align}
\end{subequations}
Then we can integrate both sides of the idfference equation to obtain,
for ideal gases,
\begin{equation}
\boxed{S - S_{0} = \int^{T}_{T_{0}}\frac{c_{V}(T')}{T'}\D T' + R\log(V/V_{0}),}
\end{equation}
where $T_{0}$ and $V_{0}$ are reference temperature and volume,
$S_{0}=S(T_{0},V_{0})$ is the reference entropy, and $\log(-)$ is the
natural logarithm.
\end{node}

\begin{node}
For perfect gases, $c_{V}$ is constant, which allows us to obtain
\begin{subequations}
  \begin{align}
S-S_{0} &= c_{V}\log(T/T_{0})+R\log(V/V_{0})\\
&= c_{V}\log(T/T_{0})-R\log(\rho_{0}/\rho),
  \end{align}
\end{subequations}
since $\rho=1/V$.
\end{node}

\begin{node}
Now, for isentropic processes, 
\begin{subequations}
\begin{equation}
S-S_{0}=0,
\end{equation}
which means perfect gases satisfy
  \begin{align}
c_{V}\log(T/T_{0})&=R\log(\rho_{0}/\rho),\\
\intertext{or equivalently,}
T/T_{0} &= (\rho/\rho_{0})^{\gamma-1},
  \end{align}
\end{subequations}
where $\gamma$ is the ratio of specific heats. We can go \emph{even further}
for perfect gases, since
\begin{subequations}
  \begin{align}
\frac{p}{p_{0}} &= \frac{RT/V}{RT_{0}/V_{0}}\\
&= \frac{T}{T_{0}}\frac{V_{0}}{V}\\
&= \left(\frac{\rho}{\rho_{0}}\right)^{\gamma},\\
\intertext{hence,}
\frac{T}{T_{0}} &= \left(\frac{p}{p_{0}}\right)^{(\gamma-1)/\gamma}.
  \end{align}
\end{subequations}
The upshot is that it suffices to know $T$, $T_{0}$, $p$, $p_{0}$ for
experiments involving isentropic processes on perfect gases.
\end{node}

\begin{node}
If $\gamma=1$, then temperature $T$ is constant for isentropic
processes (of perfect gases). In other words, isentropic processes are
isothermal for $\gamma=1$.
\end{node}

\begin{node}[Real gases]
Everything derived \emph{before} the ideal gas assumption~(\S\ref{assume:fall:lecture17:start-of-ideal-gases}) holds in general.
For real gases, there is a critical temperature $T_{c}$
which governs the phase behaviour of the gases.
\begin{center}
\includegraphics{img/img.34}
\end{center}
When the isotherm occurs at temperature $T\gg T_{c}$, then it
resembles an ideal gas. When the temperature is below the critical
threshold $T<T_{c}$ (the red isotherm), then its behaviour varies with
discontinuities at the points. For example, when $V<V_{1}$ the red
isotherm is in a liquid phase, when $V_{1}<V<V_{2}$ it is in a liquid
and gas phase, and when $V_{2}<V$ it is in a gas phase.
\end{node}

\begin{node}
The compressibility factor for a real gas is defined to be
\begin{equation}
Z := \frac{(pV)_{\text{real}}}{(pV)_{\text{ideal}}}=\frac{pV}{RT}.
\end{equation}
This is $Z=1$ for ideal gases, and $Z\neq1$ describes deviatory
behaviour.

For realistic gases, we can plot compressibility factor $Z$ against $p/p_{c}$:
\begin{center}
\includegraphics{img/img.35}
\end{center}
The reference $Z=1$ line is plotted for ideal gas behaviour (which
should hold when $T/T_{c}\gg1$). The farther we deviate from ideal
behaviour, the more important intermolecular attraction becomes.
\end{node}

\begin{node}
There are more realistic models for real gases, like the van der Waals
equation of state
\begin{equation}
p = \frac{RT}{V-b}-\frac{a}{V^{2}},
\end{equation}
which recovers the ideal gas Law when $a=b=0$. We try to empirically
fit the parameters $a$ and $b$ to lab results (intuitively, $b$
describes something like the ``excluded'' molar volume, and $a$ describes the
intermolecular force strength). There are many others, but we will not
worry about them.
\end{node}

\begin{node}[Recap]
Now, for incompressible flows, the continuity equation and the
conservation of linear momentum constituted a closed system of
equations we could solve. For compressible flows, we had more unknowns
than equations, which is why we turned to Thermodynamics for things
like conservation of energy and temperature via ideal gas Laws (to find more
equations to close the system of equations) which allowed us to
express things like pressure as a function of temperature $p=p(T)$.
\end{node}

\begin{node}
A lot of times in practice, we assume ideal gases to close the system
of equations.
\end{node}

\subsection{Quasi One Dimensional Flows}

\begin{node}
We have already examined incompressible 1-dimensional flows, which
essentially studied liquids flowing through a cylinder. The fact that
it was through a cylinder (i.e., with constant cross sectional area)
simplified calculations considerably. But what if the cross sectional
area varies along the pipe? That's what \define{Quasi One-Dimensional Flows}
studies.
\end{node}

\begin{assume}
There are two assumptions to quasi 1-dimensional flows:
\begin{enumerate}
\item Flow is \emph{mainly} in one direction, which we choose to be
  the $x$ direction. This means that $u\gg v$ and $u\gg w$.
\item Change of flow in the $x$-direction is much greater than the
  change of flow in the $y$-direction (and the $z$-direction).
\end{enumerate}
\end{assume}

\begin{example}
Consider the following nozzle:
\begin{center}
\includegraphics{img/img.36}
\end{center}
We assume the flow into the nozzle has initial velocity  $u_{1}$,
initial pressure $p_{1}$, and initial density $\rho_{1}$ (and the
initial cross sectional area is $A_{1}$). The flow out of the nozzle
has velocity $u_{2}$, pressure $p_{2}$, density $\rho_{2}$, and the
cross-sectional area is $A_{2}$.

We draw the control volume in dashed lines; although we \emph{draw}
the control volume ``slightly away'' from the surface, this is just
for clarity---in fact, it \emph{includes} the control surface. We
doodle it:
\begin{center}
\includegraphics{img/img.37}
\end{center}
We will write down the equations of motion in integral form.
Recall Eq~\eqref{eq:fall:lec3:mass-conservation} the continuity equation (in integral form):
\begin{equation}
\partial_{t}\iiint_{V}\rho\,\D V = \rho_{1}u_{1}A_{1}+\rho_{2}u_{2}A_{2}=0.
\end{equation}
The equations of motion [``conservation of linear momentum'' (\S\S\ref{node:fall:lec04:conservation-of-linear-momentum}~\textit{et~seq.})] in the
$x$-direction:
\begin{subequations}
\begin{equation}
\partial_{t}\iiint_{V}\rho u\,\D V = p_{1}A_{1}-p_{2}A_{2} + \int_{\text{walls}}p\,\D A + \int F_{x}\,\D S
\end{equation}
where we will talk about that $\int\dots\D S$ term in a moment (it is
because the flow is not in a cylinder with constant cross sectional area).
We saw in one of the first lectures how to use the Reynold's Transport
Theorem to rewrite the left-hand side as:
\begin{equation}
\rho_{1}u_{1}^{2}A_{1}+\rho_{2}u_{2}^{2}A_{2} = p_{1}A_{1}-p_{2}A_{2} + \int_{\text{walls}}p\,\D A + \int F_{x}\,\D S
\end{equation}
\end{subequations}
Then we had the conservation of energy (\S\ref{node:fall:lec04:conservation-of-energy}) equation
\begin{subequations}
\begin{equation}
\partial_{t}\int_{V}\rho(e+\frac{1}{2}u^{2})\D V = p_{1}u_{1}A_{1}-p_{2}u_{2}A_{2}-\int\dot{Q}\,\D S
\end{equation}
where we can simplify the left-hand side to
\begin{equation}
  -\rho_{1}(e_{1}+\frac{1}{2}u_{1}^{2})u_{1}A_{1}
  +\rho_{2}(e_{2}+\frac{1}{2}u_{2}^{2})u_{2}A_{2} =
  \underbrace{p_{1}u_{1}A_{1}-p_{2}u_{2}A_{2}}_{\text{work done by pressure}}-\int\dot{Q}\,\D S.
\end{equation}
\end{subequations}
The $-\int\dot{Q}\,\D S$ describes the heat transfer to surrounding,
which means $-\dot{Q}$ is energy loss per unit time per unit area.
\end{example}

\begin{node}
Now, in our diagram, we consider the cross sectional region of width $\D x$
as drawn in red:
\begin{center}
\includegraphics{img/img.38}
\end{center}
We ``zoom in'' on the cross sectional region, which should look like:
\begin{center}
\includegraphics{img/img.39}
\end{center}
The perimeter $P(x)$ is drawn in dashed red, the width of the cross
sectional area remains $\D x$. We can approximate the pipe (at this
scale) as an infinitesimally tall cylinder. Then we have
\begin{equation}
\D S = P(x)\,\D x.
\end{equation}
In this setting, we make the following approximations:
\begin{subequations}
  \begin{align}
    \int f\,\D V &\mapsto f(x)A\,\D x,
    \intertext{i.e., $f$ is approximately constant and $A$ is the
      cross sectional area,}
    \partial_{t}\int f\,\D V &\mapsto\partial_{t}(fA)\,\D x\\
    f_{2}-f_{1} &\mapsto \D f\\
    \int f\,\D A &\mapsto f\frac{\D A}{\D x}\,\D x\\
    \int f\,\D S &\mapsto fP(x)\,\D x.
  \end{align}
\end{subequations}
These approximations simplify the equations of motion.

The continuity equation becomes
\begin{equation}
\partial_{t}(\rho A)+\partial_{x}(\rho uA)=0.
\end{equation}
But $A=A(x)$, so we can write this as:
\begin{equation}
A\partial_{t}(\rho)+\partial_{x}(\rho uA)=0.
\end{equation}
Great.

Now, the conservation of $x$-momentum becomes,
\begin{equation}
\partial_{t}(\rho u A) + \partial_{x}(\rho u^{2}A) + \partial_{x}(pA)
- p\partial_{x}A=PF_{x},
\end{equation}
where we move the pressure terms to the left-hand side.

The conservation of energy simplifies to:
\begin{equation}
\partial_{t}[\rho(e+\frac{1}{2}u^{2})A] + \partial_{x}[\rho(e+\frac{1}{2}u^{2})Au]
+\partial_{x}(\rho u A) = -\dot{Q}P.
\end{equation}
In a way, this is very similar to shrinking the control volume in the
process to obtain the differential form of the Navier--Stokes
equations, but here we are less rigorous.
\end{node}

\begin{node}
We can go farther, and write the equations of motion for this
infinitesimal cylinder in a somewhat more pleasant form. The
continuity equation becomes:
\begin{equation}
\partial_{t}\rho + \partial_{x}(\rho u) = -\frac{\rho u}{A}\frac{\D A}{\D x}
\end{equation}
which tells us how the mass changes in terms of the area of the
``infinitesimally long pipe''.

The conservation of momentum,
\begin{equation}
\frac{Du}{Dt}=\partial_{t}u+u\partial_{x}u=\frac{-1}{\rho}\partial_{x}p+f
\end{equation}
where
\begin{equation}
f = \frac{PF_{x}}{\rho A}
\end{equation}
is the friction force.

The conservation of energy, using total enthalpy $h_{0}=e+\frac{1}{2}u^{2}+(p/\rho)$,
\begin{equation}
\frac{Dh_{0}}{Dt}=\partial_{t}h_{0}+u\partial_{x}h_{0}=\frac{1}{\rho}\partial_{t}p-\frac{\dot{Q}P}{\rho A}.
\end{equation}
\end{node}

\begin{node}
OK, does this make sense? Well, look at the special case when we have
a cylinder:
\begin{equation}
\frac{\D A}{\D x}=0.
\end{equation}
If friction vanishes $f=0$ and there is no heat transfer $\dot{Q}=0$,
then we recover the 1-dimensional Euler equations with $u=u(x,t)$ and $v=w=0$.
So, that's a good ``evaluation'' check.
\end{node}

\subsection{Steady Quasi One-Dimensional Flow}

\begin{node}
The first thing we should do when we obtain the equations of motion
for a flow: look at the steady flow. Then the continuity equation
becomes
\begin{equation}
\frac{\D}{\D x}(\rho u)=-\frac{\rho u}{A}\frac{\D A}{\D x},
\end{equation}
or equivalently
\begin{equation}
\frac{\D(\rho u A)}{\D x}=0.
\end{equation}
The conservation of $x$-momentum becomes
\begin{equation}
u\partial_{x}u+\frac{1}{\rho}\partial_{x}p=f.
\end{equation}
The conservation of energy becomes
\begin{equation}
\partial_{x}h+u\partial_{x}u=-\dot{q}
\end{equation}
where
\begin{equation}
h=e+\frac{p}{\rho},\quad\mbox{and}\quad\dot{q}=\frac{\dot{Q}P}{\rho A}.
\end{equation}
\end{node}

\begin{node}
When $f=\dot{q}=0$, we find the equations of motion simplify further to:
\begin{equation}
\frac{\D}{\D x}(\rho u A)=0,
\end{equation}
\begin{equation}\label{eq:lec18:fall:momentum-conservation-for-steady-quasi-1d-flow}
u\partial_{x}u + \frac{1}{\rho}\partial_{x}p=0
\end{equation}
\begin{equation}\label{eq:lec18:fall:energy-conservation-for-steady-quasi-1d-flow}
\partial_{x}h+u\partial_{x}u=0.
\end{equation}
This uses the fact that
\begin{equation}
\frac{\D h_{0}}{\D x}=\frac{\D h}{\D x} +
\frac{\D(\frac{1}{2}u^{2})}{\D x}=\frac{\D h}{\D x}+u\frac{\D u}{\D x}.
\end{equation}
\end{node}

\begin{node}
Now, returning to our equation for entropy,
\begin{subequations}
  \begin{align}
T\,\D S &= \D h - \frac{1}{\rho}\D p\\
&= -u\,\D u - \frac{1}{\rho}\D p\quad\mbox{by Eq~\eqref{eq:lec18:fall:energy-conservation-for-steady-quasi-1d-flow}}\\
&= -u\,\D u - \frac{1}{\rho}(-\rho u\,\D u)\quad\mbox{by Eq~\eqref{eq:lec18:fall:momentum-conservation-for-steady-quasi-1d-flow}}\\
&= 0.
  \end{align}
\end{subequations}
That is to say, we have an isentropic process in which the flow has a
constant entropy line.
\end{node}